% sections/en/6-conclusion.tex -- summary, limitations (one per boundary row), outlook.
\section{Conclusions}\label{sec:conclusion}
This paper decomposed the fluid gap of the multiplicity flexible job-shop scheduling problem exactly into a rounding loss
and a sequencing residual, and built an evolutionary search around that identity: a quota chromosome whose dimension is
fixed by the type structure, a ranker scoped to the residual with a bound-preserving output, and a safeguarded preselection
under exact decoding. \pending{On the fixed design set, the public benchmark and the extrapolation tier, the method reached
gaps no worse than the strongest baseline in each family within the same wall clock, with fewer exact decodings, and the
rounding gap followed the sign predicted by the theory.}

The study has the following limitations, each tied to a boundary of the design. The instance family fixes $R = 8$ job types
and $m = 10$ machines with $\bar N \le 50$; extension in the number of types is not evaluated. Setup times, transportation,
breakdowns, dynamic arrivals and lot splitting are excluded (assumptions A4 and A5). The objective is the makespan only. The
data are the parameter-table distribution and the Brandimarte benchmark, not factory data, and the Hurink sets are not
included. Commercial solvers under long time limits and algorithms designed automatically by large language models are not
compared. \pending{The learning-based baselines were reimplemented where official code was unavailable, and the
reimplementation protocol is stated with the results.}

Three directions follow. First, the same decomposition applies to any relaxation that is tight in aggregate, which suggests
testing it on batch and lot-streaming variants. Second, the ranker can be trained across type structures by sharing the
quotient-graph encoder. Third, the safeguarded union can carry other candidate sources, such as exact small-neighbourhood
searches, without changing the acceptance rule.
